\chapter{The Middleware: Pool, Blob Channel and Connector}
\label{ch:blobs}

The previous chapter described what an RDMA interface offers to a program:
endpoints with their own send queues, registered memory, one-sided writes,
and completions whose order is guaranteed only within one connection. The
middleware is built on those primitives from the bottom up, as
each layer exists to remove one cost identified in the problem statement. The
dual-endpoint pool removes the shared queue of \cref{fig:hol} by giving each
traffic class its own endpoint and its own completion handling, the blob
channel turns raw writes into named byte strings of any length under flow
control, and the connector attaches the channel to the serving engine, so
that the cache of a real request crosses the link. The components that choose
a class, an encoding and a destination for each transfer are the subject of
the next chapter, and they appear here only at the points where a decision
enters the path.

\section{One Transfer Through the Middleware}
\label{sec:overview}

\Cref{fig:architecture} shows the layers and the path that one layer of cache
takes through them. The serving engine calls the connector once per attention
layer during prefill, the connector gathers the blocks of that layer and
copies them to host memory, and the codec frames them as one blob, either as
raw 16-bit values or in the packed format of the quantizer
(\cref{ch:shrink}). The sender then obtains a class for the blob from the
classifier, and the blob channel cuts the blob into chunks, takes one credit
per chunk, and issues one RDMA write with immediate data per chunk on the
endpoint of that class. Below the channel, the pool posts each operation
through libfabric or through raw verbs, which are two implementations of one
C interface, and the network card or the Soft-RoCE driver carries the
packets. On the decode server the same layers run in the opposite order, and
the only traffic in the reverse direction is the 8-byte credit that frees a
receive slot, which always travels on the hot endpoint.

\begin{figure}[htbp]
  \centering
  \begin{tikzpicture}[diagram,
      alt={Layer diagram of the middleware as two columns joined at the bottom. The left
        column is the prefill server and the right column is the decode server. Left column,
        top to bottom, with arrows pointing down: serving engine, prefill pass, called per
        layer; connector, gathers blocks and copies them to host memory; codec, raw or packed
        by the quantizer; sender, class from the static classifier; blob channel, chunks,
        credits and staging slots; then a framed pair, the hot endpoint of the pool with depth
        64 and busy polling, and the cold endpoint of the pool with depth 512 and a blocking
        wait. The right column holds the same layers with arrows pointing up: the two
        endpoints of the pool; blob channel, reassembly and return of credits; receiver,
        staging area by key; codec, unpack and reconstruct; connector, writes blocks into the
        block table; serving engine, decode steps. Three horizontal arrows join the endpoints
        of the two servers: a solid arrow from the left hot endpoint to the right hot endpoint
        labeled marker, a dashed arrow from the right hot endpoint back to the left hot
        endpoint labeled credits, and a solid arrow from the left cold endpoint to the right
        cold endpoint labeled chunks. A wide box under both columns reads: libfabric or raw
        verbs on one RDMA device and one link, Soft-RoCE on commodity Ethernet or RoCE version
        2 on 100 gigabit Ethernet. Below it a dashed box reads: decision components of the
        next chapter, evaluated separately from this path: learned
        classifier, rule tree and controller, which choose class, compression and spill, and
        the storage tier on an NVMe drive. A dashed arrow runs from this box to the sender.},
      lay/.style={tint=#1, text width=54mm, minimum height=11mm, font=\footnotesize, align=center, inner xsep=3pt, inner ysep=3pt}]
    % prefill server
    \node[heading] at (0,1.0) {Prefill server};
    \node[lay=dGray]   (pe) at (0,0)      {Serving engine\\prefill pass, called per layer};
    \node[lay=dBlue]   (pc) at (0,-1.45)  {Connector\\gather blocks, copy to host memory};
    \node[lay=dBlue]   (pk) at (0,-2.9)   {Codec\\raw, or packed by the quantizer};
    \node[lay=dPurple] (ps) at (0,-4.35)  {Sender\\class from the static classifier};
    \node[lay=dGreen]  (pb) at (0,-5.8)   {Blob channel\\chunks, credits, staging slots};
    \node[lay=dOrange] (ph) at (0,-7.45)  {Pool, hot endpoint\\depth 64, busy polling};
    \node[lay=dOrange] (pl) at (0,-8.8)   {Pool, cold endpoint\\depth 512, blocking wait};
    \begin{scope}[on background layer]
      \node[frame, fit=(ph)(pl), inner sep=5pt] (pp) {};
    \end{scope}
    % decode server
    \node[heading] at (8.6,1.0) {Decode server};
    \node[lay=dGray]   (de) at (8.6,0)      {Serving engine\\decode steps};
    \node[lay=dBlue]   (dc) at (8.6,-1.45)  {Connector\\write blocks into the block table};
    \node[lay=dBlue]   (dk) at (8.6,-2.9)   {Codec\\unpack, reconstruct};
    \node[lay=dPurple] (ds) at (8.6,-4.35)  {Receiver\\staging area, by key};
    \node[lay=dGreen]  (db) at (8.6,-5.8)   {Blob channel\\reassembly, return of credits};
    \node[lay=dOrange] (dh) at (8.6,-7.45)  {Pool, hot endpoint\\depth 64, busy polling};
    \node[lay=dOrange] (dl) at (8.6,-8.8)   {Pool, cold endpoint\\depth 512, blocking wait};
    \begin{scope}[on background layer]
      \node[frame, fit=(dh)(dl), inner sep=5pt] (dp) {};
    \end{scope}
    % vertical flow
    \draw[flow] (pe) -- (pc); \draw[flow] (pc) -- (pk); \draw[flow] (pk) -- (ps);
    \draw[flow] (ps) -- (pb); \draw[flow] (pb) -- (pp);
    \draw[flow] (dp) -- (db); \draw[flow] (db) -- (ds); \draw[flow] (ds) -- (dk);
    \draw[flow] (dk) -- (dc); \draw[flow] (dc) -- (de);
    % logical connections between endpoints
    \draw[flow]  ($(ph.east)+(0,0.22)$) -- ($(dh.west)+(0,0.22)$) node[midway, above=1pt, small text] {marker};
    \draw[reply] ($(dh.west)+(0,-0.22)$) -- ($(ph.east)+(0,-0.22)$) node[midway, below=1pt, small text] {credits};
    \draw[flow, line width=1.6pt] (pl.east) -- (dl.west) node[midway, above=1pt, small text] {chunks};
    % device and link
    \node[tint=dGray, text width=140mm, minimum height=11mm, font=\footnotesize, align=center, inner xsep=3pt, inner ysep=3pt]
      (link) at (4.3,-10.45) {libfabric or raw verbs on one RDMA device and one link:\\Soft-RoCE on commodity Ethernet, or RoCEv2 on 100 gigabit Ethernet};
    \draw[flow] (pp.south) -- (pp.south |- link.north);
    \draw[flow] (dp.south |- link.north) -- (dp.south);
    % decision components
    \node[draw=dInk, dashed, line width=0.8pt, rounded corners=3pt, fill=white, text width=140mm,
          minimum height=15mm, font=\footnotesize, align=center, inner xsep=3pt, inner ysep=4pt]
      (dec) at (4.3,-12.35) {Decision components of the next chapter, evaluated separately from this path:\\learned classifier, rule tree and controller (class, compression, spill),\\and the storage tier (NVMe drive)};
    \draw[reply] (dec.west) -- ++(-0.45,0) |- (ps.west);
  \end{tikzpicture}
  \caption[The layers of the middleware and the path of one layer of cache]{The layers of the middleware and the path of one layer of cache. Horizontal arrows are logical connections between endpoints; both classes use the same device and cable.}
  \label{fig:architecture}
\end{figure}

\Cref{alg:transfer} states the same path as a procedure, with a reference on
each line to the algorithm that expands it. The procedure has four actors:
the engine thread of each server, which calls the connector, a sender thread
on the prefill server, and a receiver thread on the decode server. The two
engine threads never wait on each other directly, as they are coupled only
through the staging area of the decode server, in which the receiver thread
stores each blob under its key until the engine asks for it. This coupling is
what allows a layer that is lost or late to cost a timeout and a
recomputation instead of an incorrect result (\cref{sec:connector}).

\begin{pseudocode}{One layer of cache of request $r$, from the prefill server to the decode server}{alg:transfer}
  \Statex \textbf{Prefill server, engine thread,} after attention layer $\ell$ of request $r$
  \State $\var{blocks} \gets$ the blocks of $r$ in the cache of layer $\ell$, gathered on the GPU
  \State $\var{blob} \gets \Call{Encode}{\var{blocks}}$ \Comment{copy to host; raw, or \cref{alg:quant}}
  \State $c \gets \Call{Classify}{\var{key}(r,\ell),\ |\var{blob}|,\ \ell}$ \Comment{hot or cold, \cref{alg:static}}
  \State append $(\var{key}(r,\ell),\ \var{blob},\ c)$ to the send queue
  \Statex \textbf{Prefill server, engine thread,} after the last layer of $r$
  \State append $(\var{key}(r,\text{done}),\ \text{marker},\ \text{hot})$ to the send queue
  \Statex \textbf{Prefill server, sender thread}
  \For{each $(k, \var{blob}, c)$ taken from the send queue, in order}
    \State \Call{SendBlob}{$k, \var{blob}, c$} \Comment{chunks and credits, \cref{alg:blob-send}}
  \EndFor
  \Statex \textbf{Decode server, receiver thread}
  \Loop
    \State $(k, \var{blob}) \gets \Call{ReceiveBlob}{\,}$ \Comment{\cref{alg:blob-deliver}}
    \State store $\var{blob}$ in the staging area under $k$
  \EndLoop
  \Statex \textbf{Decode server, engine thread,} before the first decode step of $r$
  \For{each layer $\ell$}
    \State $\var{blob} \gets$ the entry under $\var{key}(r,\ell)$ in the staging area, or none after the timeout
    \If{no entry arrived}
      \State report the blocks of layer $\ell$ as failed \Comment{the engine recomputes them}
    \Else
      \State $\var{blocks} \gets \Call{Decode}{\var{blob}}$
      \State write $\var{blocks}$, except the prefix already held, into the block table on the GPU
    \EndIf
  \EndFor
\end{pseudocode}

Two statements about \cref{fig:architecture} prevent a misreading of the
results. First, the solid path is the one that the connector executes: the
class of a blob is chosen by the static classifier (\cref{alg:static}), the
encoding is fixed by configuration for a whole run, and no blob is diverted
to storage. On this path every layer blob carries a new key, so the access
counts of the static classifier never exceed one and the class is decided by
size alone. The marker and the credits therefore use the hot endpoint, and
every layer of cache, which is at least one 16-token block of 32 KiB for the
model used here (\cref{eq:kvsize}), uses the cold endpoint. Second, the
learned classifier, the rule tree, the controller and the storage tier are
implemented and tested as separate components, and each has a defined entry
point into the path: the sender accepts a class that overrides the
classifier, and the storage tier accepts a block together with a decision.
Connecting those entry points to the connector is future work. The experiments are assigned accordingly, as E1 to E5
exercise the pool, E6 the blob channel, and E7 and E8 the connector on the
solid path, whereas the decision components, the quantizer and the storage
tier are evaluated separately, on synthetic traffic, on caches taken from the
model, and on an NVMe drive.

\fixnote{Confirm against the connector at the commit used for the runs: the two statements above, the single send queue of the sender, and the marker that the decode side does not wait for (\cref{sec:connector}). If the decision components are connected to the serving path first, rewrite this paragraph and the dashed box of \cref{fig:architecture}.}

\section{The Dual-Endpoint Pool}
\label{sec:pool}

The pool is the lowest layer of the middleware and the one that carries its
central idea, which is that urgent and bulk traffic do not share a queue. It
sits directly on the interface of \cref{ch:libfabric}, and every layer above
it, from the blob channel to the connector, sends through it.

\subsection{One Interface, Two Libraries}

The interface of the pool is a small C header (\cref{app:calls}). It opens a
pool with a hot and a cold class, connects or listens, exposes one registered
memory region per side, and offers write, read and send per class, together
with calls that poll for arrivals and wait for local completions. Two
libraries implement this header: libscuffedfabric on libfabric, and
libscuffedverbs on raw verbs with reliable connected queue pairs. The
benchmark of \cref{sec:sfbench} is one source file linked against either
library, so that when the two differ in a measurement, the library beneath
the header is the only element that changed, which is the comparison that RQ2
requires.

\subsection{Two Classes, Two Queues}

The pool gives each traffic class its own endpoint and its own completion
handling, and both halves of that statement matter. The separate endpoint
removes the shared send queue of \cref{fig:hol}, so that the wait of
\cref{eq:fifo} no longer contains cold messages. The separate completion
queue removes a second and less visible queue: if both classes reported to
one completion queue, a hot completion that arrived behind $k$ cold
completions would cost $k$ polls to reach, whereas with a queue per class the
hot poller never sees a cold completion. The two classes are then configured
for opposite goals, as \cref{tab:classes} lists.

\begin{table}[htbp]
  \centering
  \caption[Configuration of the two traffic classes]{Configuration of the two traffic classes.}
  \label{tab:classes}
  \small
  \begin{tabular}{L{3.6cm}L{5.2cm}L{5.2cm}}
    \toprule
    & Hot class & Cold class \\
    \midrule
    Goal & lowest latency for small messages & highest throughput for large ones \\
    Queue depth & 64 & 512 \\
    Waiting for completions & busy polling, or a counter & blocking wait, so no processor core spins \\
    Small payloads & copied by the provider at the call, so the buffer is free at once & not applicable \\
    Transport & reliable connection & reliable connection \\
    \bottomrule
  \end{tabular}
\end{table}

The hot class spends a processor core to save microseconds, as its thread
polls the completion queue without sleeping and small payloads are copied by
the provider at the call, which frees the buffer at once. The cold class does
the opposite, because a thread that streams megabytes should sleep while the
link is busy, and a cold thread that polled would compete with the hot poller
for the processor. For this reason the cold class waits on a kernel wait
object, and the cluster runs pin the two threads to different cores
(\cref{sec:sfbench}).

\Cref{alg:pool} gives the data path of the pool in full. An operation is
submitted to the endpoint of its class, and if the send queue of that
endpoint is full, the caller reaps completions of the same class until space
appears, so that a full cold queue delays cold callers only. Completions are
counted per class, by busy polling on the hot class and by a blocking wait of
at most 10 ms on the cold class.

\begin{pseudocode}{Posting an operation on a class and counting its completion}{alg:pool}
  \Function{Post}{$c$, operation} \Comment{$c$ is hot or cold}
    \Loop
      \State submit the operation to the endpoint of class $c$ \Comment{write, read or send}
      \If{the provider accepted it}
        \State $\var{posted}[c] \gets \var{posted}[c] + 1$ and \Return
      \ElsIf{the error is not a full queue, or the queue has been full for 60 s}
        \State \textbf{fail} with that error
      \EndIf
      \State \Call{Progress}{$c$} \Comment{reap completions to free queue space}
    \EndLoop
  \EndFunction
  \Function{Progress}{$c$}
    \State read up to 64 entries from the completion queue of class $c$
    \If{an entry reports an error}
      \State \textbf{fail} with that error \Comment{never counted as done}
    \EndIf
    \State $\var{done}[c] \gets \var{done}[c] +{}$ the number of entries read
  \EndFunction
  \Function{WaitTx}{$c$, $n$} \Comment{returns once $\var{done}[c] \ge n$}
    \While{$\var{done}[c] < n$}
      \If{class $c$ has a wait object} \Comment{cold class}
        \State sleep on its completion queue for at most 10 ms, then count as in \textsc{Progress}
      \Else \Comment{hot class}
        \State \Call{Progress}{$c$}
      \EndIf
    \EndWhile
  \EndFunction
\end{pseudocode}

Two details of \cref{alg:pool} carry weight in later sections. An entry with
an error status ends the operation with that error and is never added to
$\var{done}$, which is rule 2 of \cref{sec:rules} enforced at the lowest
layer. In addition, the counters $\var{posted}[c]$ and $\var{done}[c]$ give
every operation on a class a sequence number, and since the completions of
one reliable connection are reported in posting order, $\var{done}[c] \ge n$
means that the $n$th operation posted on class $c$ has completed, which is
the fact on which \cref{lem:staging} rests. When the hot class is configured
with a completion counter instead of a queue, \textsc{Progress} reads the
counter into $\var{done}[c]$ and the rest is unchanged.

What the pool guarantees is narrow and is stated exactly, since the
experiments test nothing more. No cold message is ahead of a hot message in
the send queue of the hot class, and no cold completion is ahead of a hot
completion in its completion queue. What happens below the two endpoints,
where their packets meet on one device and one cable, is outside the control
of the pool: on a hardware card the device arbitrates between queue pairs,
and on Soft-RoCE the kernel does. Whether the hot class stays fast there is
the question that E3 measures (\cref{sec:res-e3}).

\subsection{Differences Between Providers}
\label{sec:providers}

The tests of the pool run over three software providers, the TCP provider,
the RXM layer on top of it, and the shared-memory provider, with libfabric
1.17.0 and with 2.7.0 built from source, and they showed three differences
that the design has to respect. The shared-memory provider does not promise
that two writes to the same peer are applied in order, so a signal that
relies on a second write arriving after the first is not valid on it. The
TCP provider did not increment the counter for remotely originated writes, so
the counter signal is skipped there. The third difference concerns progress,
the work that moves a transfer forward after the posting call has returned.
The RXM layer and the shared-memory provider report manual progress and the
TCP provider automatic progress, and in these tests a queued transfer
advanced on all three only during calls into the library
\citep{libfabric_repo}. The pool has no thread of its own, so a caller keeps
calling \textsc{Progress} or \textsc{WaitTx} until its operation is complete,
on both sides, and the blob channel is built on that rule
(\cref{sec:channel}). The RXM layer and the shared-memory provider also need
the progress of the peer for the first use of an endpoint, so the pool
exchanges one message on each such endpoint when the connection is opened.
The verbs provider, by its source, offers no counters and no events for
remotely originated writes in either release, so the counter signal of
\cref{tab:signals} cannot run on it, and a counter for the hot completions
exists there on reliable-datagram endpoints, through the RXM layer.

\waiting{results from E1 to E5 on the verbs provider: state whether the
ordering difference occurs there}

\section{The Blob Channel}
\label{sec:channel}

The pool moves bytes between offsets of two registered regions and reports
completions, which is too low a level for a caller that holds a cache layer
of several megabytes under a name. The blob channel supplies the missing
parts: a framing that names each byte string, a division into chunks that fit
a fixed slot, and a credit scheme that prevents the sender from overwriting
data that the receiver has not yet copied out.

\subsection{Layout}

Each side registers one memory region and divides it into three parts: $N$
receive slots of $B$ bytes into which the peer writes, $N$ staging slots in
which outgoing chunks are assembled, and a small scratch area for credits.
The defaults are $N = 8$ and $B = 4$ MiB, and every slot begins with a
64-byte header (\cref{fig:slot}, with the fields in \cref{app:wire}).

\begin{figure}[htbp]
  \centering
  \begin{tikzpicture}[diagram,
      alt={Byte layout of one slot of the blob channel, drawn as a horizontal bar in three
        parts. Part one, 64 bytes, the header: magic number, blob identifier, key length, flags,
        total blob length, chunk offset and chunk length. Part two, the key, a variable number
        of bytes. Part three, the chunk payload, which fills the rest of the slot up to B bytes.
        A brace under the whole bar is labeled one slot, B bytes, 4 mebibytes by default.},
      cell/.style={draw=dInk, line width=0.8pt, minimum height=11mm, font=\footnotesize, align=center, inner sep=3pt}]
    \node[cell, fill=dBlue!12, minimum width=44mm] (h) at (0,0) {Header, 64 bytes};
    \node[cell, fill=dGreen!12, minimum width=20mm, anchor=west] (k) at (h.east) {Key};
    \node[cell, fill=dOrange!10, minimum width=58mm, anchor=west] (p) at (k.east) {Chunk payload};
    \draw[decorate, decoration={brace, amplitude=5pt, mirror}, draw=dGray, line width=0.6pt]
      ($(h.south west)+(0,-0.15)$) -- ($(p.south east)+(0,-0.15)$)
      node[midway, below=7pt, small text] {one slot: $B$ bytes (4 MiB by default)};
    \node[small text, align=left, anchor=north west, text width=120mm] at ($(h.south west)+(0,-1.0)$)
      {Header fields: magic number, blob identifier, key length, flags, total length, chunk offset, chunk length.};
  \end{tikzpicture}
  \caption[One slot of the blob channel]{One slot of the blob channel.}
  \label{fig:slot}
\end{figure}

\subsection{Protocol}

A blob is a key and a byte string, and a credit is the right to write into
one specific receive slot of the peer, of which the sender initially holds
all $N$. \Cref{alg:blob-send,alg:take-credit} give the sending side. The
sender cuts the blob into chunks of at most $B - 64 - |\var{key}|$ bytes,
which is what remains of a slot of $B$ bytes after the 64-byte header and the
key of $|\var{key}|$ bytes, and for each chunk it takes a credit, builds header, key and chunk in the staging slot with the
same number, and writes that staging slot into the receive slot of the peer
with one RDMA write with immediate data, in which the immediate value carries
the kind ``data'' and the slot number.

\begin{pseudocode}{Sending a blob}{alg:blob-send}
  \Function{SendBlob}{$\var{key}$, $\var{data}$, class $c$}
    \State $\var{cap} \gets B - 64 - |\var{key}|$ \Comment{payload bytes per chunk}
    \State $\var{id} \gets$ next blob identifier, $\var{off} \gets 0$
    \Repeat
      \State $\var{len} \gets \min(\var{cap},\ |\var{data}| - \var{off})$
      \State $s \gets \Call{TakeCredit}{\,}$ \Comment{a free receive slot of the peer, \cref{alg:take-credit}}
      \State build header, $\var{key}$ and $\var{data}[\var{off} : \var{off}+\var{len}]$ in staging slot $s$
      \State \Call{Post}{$c$, write of staging slot $s$ to receive slot $s$, immediate (data, $s$)}
      \State $\var{mark}[s] \gets (c,\ \var{posted}[c])$
      \State $\var{off} \gets \var{off} + \var{len}$
    \Until{$\var{off} \ge |\var{data}|$}
    \State \Call{Drain}{\,} until $\var{done}[c] \ge \var{posted}[c]$ \Comment{every chunk has completed at the sender}
  \EndFunction
\end{pseudocode}

\begin{pseudocode}{Taking a credit before a chunk is written}{alg:take-credit}
  \Function{TakeCredit}{\,}
    \Loop
      \If{a credit $s$ is free}
        \If{$\var{mark}[s]$ is set to $(c', n)$}
          \State \Call{Drain}{\,} until $\var{done}[c'] \ge n$ \Comment{its earlier write is complete}
        \EndIf
        \State \Return $s$
      \EndIf
      \State \Call{Drain}{\,} \Comment{\cref{alg:blob-recv}; fail once the deadline has passed}
    \EndLoop
  \EndFunction
\end{pseudocode}

\Cref{alg:blob-recv,alg:blob-deliver} give the receiving side. The first of
them is also the routine that a waiting sender runs, since credits arrive
through the same completion queues as data. The completion of each write names the slot through its
immediate value. The receiver copies the chunk out of the slot to its place
in the blob, by blob identifier and offset, and only then returns the credit
of the slot, as an 8-byte write on the hot endpoint whose immediate value
carries the kind ``credit'' and the slot number. When all bytes of a blob
have been copied, the blob is delivered to the caller under its key.

\begin{pseudocode}{Draining both endpoints}{alg:blob-recv}
  \Function{Drain}{\,}
    \For{class $c$ in (hot, cold)}
      \State \Call{Progress}{$c$}
      \For{each arrival on the endpoint of $c$ with immediate $(\var{kind}, s)$}
        \If{$\var{kind}$ is credit}
          \State add $s$ to the free credits
        \Else \Comment{a data chunk is in receive slot $s$}
          \State read $\var{id}$, $\var{key}$, $\var{total}$, $\var{off}$ and $\var{len}$ from the header of slot $s$
          \State copy $\var{len}$ payload bytes of slot $s$ to offset $\var{off}$ of the buffer of blob $\var{id}$
          \State \Call{Post}{hot, 8-byte write with immediate (credit, $s$)}
          \If{all $\var{total}$ bytes of blob $\var{id}$ have been copied}
            \State append $(\var{key},\ \text{buffer})$ to the ready list
          \EndIf
        \EndIf
      \EndFor
    \EndFor
  \EndFunction
\end{pseudocode}

\begin{pseudocode}{Receiving a blob}{alg:blob-deliver}
  \Function{ReceiveBlob}{\,}
    \While{the ready list is empty}
      \State \Call{Drain}{\,} \Comment{\cref{alg:blob-recv}}
    \EndWhile
    \State \Return the first entry of the ready list
  \EndFunction
\end{pseudocode}

Three choices in these algorithms carry the properties proved below. A
chunk and its notification are a single operation, the write with immediate
data, so that no second message exists that could overtake the data. The
receiver posts the credit only after it has copied the chunk out, so that a
slot is never rewritten while it is being read. Finally, the credit travels
on the hot endpoint whatever class the data used, so that the signal which
allows the sender to continue never waits behind bulk data that the same side
is sending in the opposite direction.

A fourth choice concerns delivery, and it was made after a test failed.
\textsc{SendBlob} returns when every chunk has completed at the sender, and
it drains both endpoints while it waits. An earlier version returned as soon
as the chunks were posted, and blobs of several chunks then stalled
intermittently on the TCP provider when the sender made no further call: by
the progress rule of \cref{sec:providers}, the part of a write that the
socket had not accepted stayed in the queue of the provider. With the wait,
the reproducer of the stall passes with both library versions on a
development machine, and the suite tests a sender that stays idle after each
blob. The time recorded for a blob now ends when its last chunk has
completed, and two blobs no longer overlap on the wire unless the caller
posts without the wait and flushes later. The sender of the connector also
polls its channel every 2 ms while it is idle, so that credits are served.

\waiting{run R5: this test on the verbs provider and on RXM over verbs}

\subsection{Proved Properties}

Four properties of this protocol are proved in \cref{sec:proofs-blob}.

\begin{proposition}[Slots are never overwritten in use]
\label{prop:slots}%
A receive slot cycles through three states, free with its credit at the
sender, being written, and written with the receiver copying out, and returns
to free only when the receiver sends the credit. The sender writes a slot only
while holding its credit, so no slot is written while its previous contents
are still being read.
\end{proposition}

\begin{proposition}[Notification never precedes data]
\label{prop:notify}%
The receiver learns of a chunk only from the completion of the write with
immediate data, and by fact 3 of \cref{sec:ordering} that completion follows
the placement of every byte of the write.
\end{proposition}

\begin{proposition}[Reassembly is independent of arrival order]
\label{prop:reassembly}%
Each chunk carries its blob identifier, the total length, its own offset and
its length, so a receiver reconstructs the blob from chunks received in any
order.
\end{proposition}

\begin{proposition}[Throughput bound]
\label{prop:credit}%
Let $N$ be the number of receive slots, $B$ the size of a slot in bytes, and
$t_{\text{credit}}$ the mean duration of the round trip of a credit, from the
moment a slot is written to the moment its credit can be used again. Then the
throughput $X$ of the channel, in bytes per second over a long period,
satisfies
\mathalt{X is at most N times B divided by t credit}{%
\begin{equation}
  X \le \frac{N B}{t_{\text{credit}}}.
  \label{eq:credit}
\end{equation}}\par % keep \par: a theorem that ends with a display stops TeX Live 2025 when the page breaks after it
\end{proposition}

The last bound sizes the window. With the defaults, $N B = 32$ MiB can be in
flight. A 10 gigabit link is the tighter limit as long as $t_{\text{credit}}$
stays under about 27 ms, and a 100 gigabit link as long as it stays under
about 2.7 ms. On the
lab pair the window is therefore not the bottleneck unless credit processing
is delayed by tens of milliseconds, which is the reason credits are kept off
the cold endpoint.

\subsection{Proofs of the Channel Properties}
\label{sec:proofs-blob}

The notation is that of \cref{sec:channel}. Each side has $N$ receive slots of
$B$ bytes, numbered by $s$. A credit for slot $s$ is the right to write into
receive slot $s$ of the peer, and the sender starts with all $N$ credits. The
counters $\var{posted}[c]$ and $\var{done}[c]$ of \cref{alg:pool} count the
operations posted and completed on the endpoint of class $c$.

\subsubsection{Slots Are Never Overwritten in Use}

\begin{proof}[Proof of \cref{prop:slots}]
For a slot $s$ consider three conditions.
\begin{enumerate}
  \item[(A)] The sender holds the credit for $s$.
  \item[(B)] A write to $s$ has been posted, and the receiver has not finished
        copying its chunk out.
  \item[(C)] The receiver has finished copying, and the sender has not yet
        processed the credit message for $s$.
\end{enumerate}
The claim is that exactly one of these holds for every slot at every moment,
and the proof is by induction over the steps of the protocol.

\emph{At the start.} The sender holds every credit, so (A) holds for every
slot and the others do not.

\emph{Each step.} Three steps touch a slot. The sender takes the credit for
$s$ and posts a write to $s$, which it can do only in (A), and the slot is
then in (B). The receiver copies the chunk out and then posts the credit
message, in that order (lines 9 and 10 of \cref{alg:blob-recv}), and the slot
moves from (B) to (C). The sender processes the credit message and holds the
credit again, and the slot moves from (C) to (A). No step applies to a slot
in any other condition, and each step leaves the slot in exactly one
condition.

\emph{Conclusion.} A write is posted to $s$ only in (A), and (A) is entered
only from (C), which is entered only when the copy has finished. Hence no
slot is written while its previous contents are still being read.
\end{proof}

\Cref{fig:credit} draws the cycle that the proof follows.

\begin{figure}[htbp]
  \centering
  \begin{tikzpicture}[diagram,
      alt={Cycle of three conditions of one receive slot, drawn as three boxes joined by
        arrows into a loop. Box A: free, the sender holds the credit. An arrow from A to B is
        labeled: the sender takes the credit and posts the write. Box B: in use, the write is
        posted and the receiver has not finished copying. An arrow from B to C is labeled: the
        receiver copies the chunk, then posts the credit. Box C: copied out, the credit is on
        its way to the sender. An arrow from C back to A is labeled: the sender processes the
        credit. A note says a write is posted to the slot only in condition A.},
      st/.style={tint=#1, text width=34mm, minimum height=13mm, font=\footnotesize, align=center, inner xsep=3pt}]
    \node[st=dGreen]  (a) at (0,0)     {(A) Free\\sender holds the credit};
    \node[st=dOrange] (b) at (6.1,0)   {(B) In use\\written, being copied};
    \node[st=dBlue]   (c) at (12.2,0)  {(C) Copied out\\credit on its way back};
    \draw[flow] (a) -- (b) node[midway, above=2pt, small text, align=center] {takes credit,\\posts write};
    \draw[flow] (b) -- (c) node[midway, above=2pt, small text, align=center] {copies, then\\posts credit};
    \draw[flow] (c.south) -- ++(0,-1.0) -| (a.south);
    \node[small text] at (6.1,-1.4) {the sender processes the credit};
    \node[small text] at (6.1,-2.2) {A write is posted to the slot only in condition (A).};
  \end{tikzpicture}
  \caption[The three conditions of a receive slot]{The three conditions of a receive slot in the proof of \cref{prop:slots}.}
  \label{fig:credit}
\end{figure}

The sending side needs the matching property for its staging slots, because a
one-sided write reads its source buffer after it is posted.

\begin{lemma}[Staging slots are not reused early]
\label[lemma]{lem:staging}%
The sender does not rewrite staging slot $s$ until the earlier write that
used it as a source has completed at the sender.
\end{lemma}

\begin{proof}
When the sender posts a write from staging slot $s$ on the endpoint of class
$c$, it records $\var{mark}[s] = (c, n)$, where $n = \var{posted}[c]$ is the
number of operations posted on that endpoint so far, this one included
(line 9 of \cref{alg:blob-send}). Before it reuses $s$ it waits until
$\var{done}[c] \ge n$ (line 5 of \cref{alg:take-credit}). Completions on one
endpoint are reported in the order in which the operations were posted, so
when the number of completions has reached $n$, the recorded write is among
them. A completed write no longer reads its source, and a write small enough
to be injected is copied when it is posted, so that its source is free at
once.
\end{proof}

The lemma depends on in-order completion on a single endpoint. That holds
for the send queue of a reliable connection, which is what both back ends
use.

\subsubsection{Notification Never Precedes Data}

\begin{proof}[Proof of \cref{prop:notify}]
The chunk and its notification are one operation, a write with immediate
data. The receiver is told of the chunk only by the completion that carries
the immediate value. By fact 3 of \cref{sec:ordering}, that completion is
generated after every byte of the write is in place. There is no second
message that could arrive first.
\end{proof}

A design that sent the data in one write and a separate ``done'' message on
another endpoint would need an ordering guarantee between endpoints, and
none exists.

\subsubsection{Reassembly}

\begin{proof}[Proof of \cref{prop:reassembly}]
The sender cuts the byte range from $0$ to the total length into consecutive
intervals that do not overlap: the offset starts at $0$ and advances by the
length of each chunk until it reaches the total. Each chunk's header carries
the blob identifier, the total, the chunk's offset and its length, and the
receiver copies the chunk to exactly that range of the buffer it keeps for
that blob identifier. The position of a chunk therefore does not depend on
when it arrives. The transport delivers each write once or reports the
connection as failed, so on a connection that has not failed every interval
arrives once, and the number of bytes received equals the total exactly when
every interval has arrived. That is the condition under which the receiver
delivers the blob.
\end{proof}

The proof assumes that every chunk arrives. A transfer in which one does
not, or in which its completion is not processed, is a fault of the
implementation or of the provider and not a gap in this argument.

\subsubsection{The Throughput Bound}

\begin{proof}[Proof of \cref{prop:credit}]
Observe the channel for a period of length $t_{\text{obs}}$. Let $k_i$ be the
number of times that slot $i$ is written in that period, and let
$K = k_1 + \dots + k_N$ be the number of all writes.

\emph{Counting round trips.} Between two consecutive writes to the same slot
lies one complete round trip of its credit, so slot $i$ completes at least
$k_i - 1$ round trips inside the period, and their durations add up to at
most $t_{\text{obs}}$. Over all $N$ slots there are therefore at least
$K - N$ complete round trips with a total duration of at most
$N t_{\text{obs}}$. Their mean duration is $t_{\text{credit}}$, hence
\mathalt{the quantity K minus N, times t credit, is at most N times t obs, so K is at most N
  plus N times t obs divided by t credit}{%
\begin{equation}
  (K - N)\, t_{\text{credit}} \le N\, t_{\text{obs}}
  \qquad\text{so}\qquad
  K \le N + \frac{N\, t_{\text{obs}}}{t_{\text{credit}}}.
\end{equation}}

\emph{From writes to bytes.} Each write carries at most $B$ bytes, header and
key included, so the bytes moved in the period are at most $B K$, and the
throughput over the period is
\mathalt{X equals the bytes moved divided by t obs, which is at most B times K over t obs,
  which is at most N times B over t obs, plus N times B over t credit}{%
\begin{equation}
  X = \frac{\text{bytes moved}}{t_{\text{obs}}}
    \le \frac{B K}{t_{\text{obs}}}
    \le \frac{N B}{t_{\text{obs}}} + \frac{N B}{t_{\text{credit}}}.
\end{equation}}
The first term vanishes as $t_{\text{obs}}$ grows, which gives
\cref{eq:credit}.
\end{proof}

\section{Five Completion Signals}
\label{sec:signals}

The channel uses the completion of a write with immediate data as its arrival
signal, which is the safe choice by \cref{prop:notify} and not the only one.
It costs a posted receive per write and a completion queue entry to process,
and for a message of a few bytes those costs are a large part of the whole
exchange. \Cref{tab:signals} lists the five mechanisms that the benchmark can
run, with the guarantee on which each one rests.

\begin{table}[htbp]
  \centering
  \caption[Five completion signals for a one-sided write]{Five completion signals for a one-sided write.}
  \label{tab:signals}
  \small
  \begin{tabular}{L{3.1cm}L{5.2cm}L{5.9cm}}
    \toprule
    Mechanism & How the receiver learns & Correct because, and when it fails \\
    \midrule
    Completion with immediate data & reads its completion queue & fact 3. Costs a posted receive per write \\
    Flag in the payload & polls a byte at the end of the buffer & nothing guarantees it for a multi-packet message (fact 2). Safe for single-packet messages \\
    Sentinel write & polls a second, tiny write posted after the data & fact 1, if the provider keeps write-after-write order. The shared-memory provider does not \\
    Remote-write counter & reads a counter the provider increments & counts arrivals without identifying them. The TCP provider did not count, and the verbs provider has no counters \\
    Write, then read & the sender reads 8 bytes back and waits for that completion & fact 1: the read is executed after the write. Signals the sender, not the receiver \\
    \bottomrule
  \end{tabular}
\end{table}

The second and third rows are the lightweight signals used in low-latency
collective libraries, where a receiver polls memory instead of a queue
\citep{shen2026sol}. The fourth is the counter of fabric-lib
\citep{licker2026fabriclib} expressed in libfabric. The fifth signals the sender instead of the receiver, which matters when the
sender must know that data is in place before it tells a third party to read
it. A production file system has reported an issue of that kind, in which a
race between its RDMA layer and the file system allowed a client to read
incorrect file data \citep{ibm_ij44629}. E5 measures all five on both pairs,
with the rule that a mechanism replaces the default only if it moves the 99th
percentile.

\section{The Connector}
\label{sec:connector}

The pool and the channel are measured by a benchmark, which answers the
first two research questions and says nothing about a real request. The
connector closes that distance, as it is the piece that lets the serving
engine of \cref{ch:inference} produce its cache on one server and consume it
on another through the blob channel.

\subsection{The Interface Offered by the Serving Engine}

vLLM supports split serving through a connector interface: a class with a
scheduler side, which decides for each request whether cache will be saved or
loaded, and a worker side, which is called around each attention layer to
save or load the blocks of that layer \citep{vllm_disagg}. The engine's own
implementation of this interface uses NIXL \citep{nixl}.
ScuffedRdmaConnector implements the same interface on the blob channel, and
it is checked against the engine source at a pinned revision and against a
released version.

\subsection{One Request from End to End}

The connector follows \cref{alg:transfer} for each layer, and the steps
around it involve the proxy and the scheduler side of both instances.

\begin{enumerate}
  \item A small proxy receives the user's request and assigns it a transfer
        identifier. It sends the request to the prefill instance with the
        output length set to one token.
  \item The scheduler side of the prefill instance records the transfer
        identifier and marks the request as a producer.
  \item As the prefill pass runs, the worker side is called for each
        attention layer. It gathers the blocks of that layer from GPU memory,
        copies them to host memory, encodes them, and queues one blob keyed
        by the transfer identifier and the layer name.
  \item When the last layer is queued, a short ``done'' marker for the
        request is queued on the hot class.
  \item The proxy forwards the original request, with the transfer
        identifier, to the decode instance.
  \item The scheduler side of the decode instance reports that all but the
        last prompt token are already computed.
  \item Its worker side waits for the blob of each layer, skips the blocks
        that it already holds in a local prefix cache, decodes the rest and
        writes them into the block table of the engine on the GPU.
  \item Decoding starts and tokens stream to the user.
\end{enumerate}

Blobs are keyed by the transfer identifier of the proxy because the engine
appends a random suffix to its own request identifiers, so that the two
instances do not agree on those. The decode instance does not wait for the
marker: it starts when the proxy forwards the request and then waits for each
layer under its key, and the marker is discarded once the layers are loaded.

\subsection{Where the Classes Separate on the Serving Path}

The sender of the connector keeps one first-in first-out queue and one thread
for both classes (lines 4 to 7 of \cref{alg:transfer}), so on the serving
path the classes separate below it, in the channel and the pool. Within one
request that order is required, as the marker closes the layers that precede
it. Across concurrent requests, however, a marker that is queued behind the
layers of another request waits for them in the sender before it reaches the
hot endpoint, which reintroduces above the pool the wait that the pool
removes. A send queue per class would carry the separation up to the
connector, and it is listed as future work. The
measurements of the separation itself, E3 and E4, are taken with the
benchmark directly on the pool and do not pass through this queue.

\subsection{Failure Behavior}

A layer that does not arrive within the timeout is reported to the engine as
a failed load for those blocks, and the engine recomputes them. A lost layer
therefore costs time and never produces incorrect output, which a test
verifies by dropping a layer on purpose. The default timeout of 30 seconds
per layer is far too long for a real deployment and is a configuration value.
If the RDMA channel cannot be opened, the connector falls back to a TCP
channel and records an event, so that a fallback is visible in the metrics
instead of silent.

\subsection{The Cost of Staging Through Host Memory}

With no GPU-direct path (\cref{ch:compat}), each layer is copied from GPU
memory to host memory before it is sent, and back on the other side. A layer
of $s$ bytes takes the time $t_{\text{copy}} = s / P$ to copy, where $P$ is
the rate of the copy in bytes per second, and the time
$t_{\text{wire}} = 8 s / C$ to cross a link of rate $C$ in bits per second
(\cref{eq:wiretime}). The ratio of one copy to the wire time of the same
bytes therefore does not depend on their number:
\mathalt{copy time divided by wire time equals s over P, divided by 8 s over C, which equals
  C divided by 8 P}{%
\begin{equation}
  \frac{t_{\text{copy}}}{t_{\text{wire}}} = \frac{s / P}{8 s / C} = \frac{C}{8P}.
  \label{eq:staging}
\end{equation}}
For an assumed $P$ of 20 GB/s this ratio is about 6 percent at 10 gigabits
per second and 62 percent at 100, which means that the copy is a minor cost
on the lab pair and a major one on the faster link. The copy rate is an
assumption, and E7 times the real copies.

\subsection{Baselines}

The comparison in E8 is against the NIXL connector of the engine on two back
ends. On UCX, NIXL is in its standard configuration. On libfabric, the NIXL
revision read for this work searches for specific vendor adapters and falls
back to TCP, so on the lab pair that row measures the libfabric path of NIXL
over TCP. A third reference is our connector on its TCP channel, which
isolates the effect of the transport from the effect of the connector.

\section{NUMA-Aware Placement of the Registered Memory}
\label{sec:numa}

A server with more than one processor socket has non-uniform memory access
(NUMA). Each socket has its own memory and its own I/O buses, the sockets are
joined by an interconnect, and Linux presents each socket with its memory as
a \emph{node}. The time to reach a page and the bandwidth of the access
depend on the distance between the node that holds the page and the node of
the processor or I/O bus that makes the access \citep{kernel_numa}. A network
card and a GPU are attached to the bus of one node, so the placement of
memory matters for a device in the same way as for a thread.

The arena of the pool (\cref{sec:pool}) is touched by two parties for every
chunk. The network card reads a staging slot or writes a receive slot, and a
thread copies the chunk into or out of that slot. Where the arena lies is
decided when its pages are first written: unless a policy says otherwise,
Linux allocates a page on the node of the processor on which the allocating
thread runs \citep{kernel_mempolicy}. The pool maps the arena and writes
zeros to all of it in the thread that opens the pool, so the arena is placed
by where that thread happened to run. Registration then pins the pages for as
long as the pool exists \citep{kernel_pin}, while the scheduler remains free
to move the thread to another node, because it does not take the location of
a task's memory into account directly \citep{kernel_numa}.

On a server with two nodes this leaves three cases. With the arena and the
copying thread on the node of the card, no access crosses the interconnect.
With the arena on the node of the card and the thread on the other node,
every copy crosses it. With the arena on the other node, every transfer of
the card crosses it, wherever the thread runs. The staging copy between GPU
memory and host memory is subject to the same rule, so the rate $P$ of
\cref{eq:staging} can also depend on the placement. If the GPU and the card
are attached to different nodes, the arena is local to only one of the two,
and which to favor is a matter for measurement.

A NUMA-aware pool removes the dependence on chance in four steps. It reads
the node of the network card, which the kernel exposes for every PCI device
and reports as $-1$ when the firmware does not supply it
\citep{kernel_sysfs_pci}. It binds the mapping of the arena to that node
before the first write, because a policy governs the pages that are
allocated after it is installed \citep{kernel_mempolicy}; in the code the
step belongs between the mapping of the arena and its zeroing. It registers
the arena only then, and it restricts the threads that copy into and out of
the slots to the processors of that node. The benchmark already pins its two
threads to given processors, and the pool and the connector do not. Each
results file would then name the node of the card and of the arena beside
the number of nodes that the hardware record holds. This binding was planned
at the start of the work and is not implemented.

The measurements of this thesis are not expected to depend on the placement.
Servers 2, 3 and 4 have one processor socket each (\cref{tab:servers}), which
normally forms a single node, and for such a server the vendor of the
ConnectX cards states that no tuning of this kind is required
\citep{nvidia_cx4_tuning}. The placement can therefore not be evaluated on
this testbed, which is a limit of it. On a server with two sockets the three
cases can be compared before the binding is written, because a memory policy
that is set for a program when it is started is inherited by it
\citep{kernel_mempolicy}: the benchmark is started with its memory bound to
one node and its two threads pinned to processors of either. The binding is
the first addition to make before the pool runs on such a server.
\waiting{the hardware record (run R0): the number of nodes on each server}

The path of a transfer is thereby fixed from the engine of one server to the
engine of the other, and the next chapter turns to the choices that can be
made for each transfer along it.
